% GENERATED FILE. Edit canonical lessons, then run npm run generate:books.
% canonical-source-hash: fnv1a64:b3c00c570f1b1197
% canonical-lessons: ES-C283-disculpe, ES-C283-permiso, ES-C283-amable, ES-C283-mil-gracias, ES-C283-no-hay-de-que

\chapter{Courtesy Past Thank You --- Excuse Me, With Permission, Kind, A Thousand Thanks}
\label{ch:cortesia}
\hlchaptermodality{283}

\begin{chapteropening}
\textbf{By the end of this chapter:} \emph{I can apologise, ask to pass, call someone kind and thank them at more than one weight in Spanish.}

Complete ES-C283-no-hay-de-que: say all five, and take the last one apart into the words you already had.
\end{chapteropening}

\section[disculpe]{disculpe — take the fault off me}

\label{lesson:ES-C283-disculpe}



\begin{quote}
Before a new run: what is an \emph{invitado}? (One who has been invited.) And what was \emph{la sonrisa} built from? (\emph{Ridere}, to laugh, turned down by \emph{sub-}.)
\end{quote}

\subsection*{What to know first}
\begin{itemize}
\raggedright
  \item perdón — the apology you already have.
  \item por favor — the politeness you already have.
\end{itemize}

\begin{sounds}
\begin{itemize}
\raggedright
  \item \emph{dis-CUL-pe}, stress in the middle. The \emph{c} before \emph{u} is a hard \emph{k}, as in \emph{casa}. $\to$ reference
\end{itemize}
\end{sounds}

\begin{cousinweb}
\textbf{disculpe} — \textquotedblleft{}excuse me.\textquotedblright{}

From \textbf{disculpar}: \emph{dis-}, which undoes, plus \textbf{culpa}, fault. To \emph{disculpar} is to take the fault off someone. Said to a person — \emph{disculpe} — it asks them to un-fault you.

Latin \emph{culpa} stands intact in English \textbf{culpable} and in \textbf{mea culpa}, my fault, and worn down in \textbf{culprit}. The \emph{dis-} is the reversing prefix of \emph{disappear} and \emph{disarm}: it does not merely deny a thing, it undoes it.

\emph{Disculpe} is the \emph{usted} form, so it is the one for a stranger. To a friend you would say \emph{disculpa} — the same choice you have been making since you met \emph{tú} and \emph{usted}.
\end{cousinweb}

\begin{grammarlens}[title={What you've built}]
You now have two apologies, and they are not the same size.

\emph{Perdón} asks for pardon, which is a thing granted. \emph{Disculpe} asks for the fault to be lifted, which is a thing undone. Spanish gives you both, and a language with two words for one moment is telling you the moment has two parts.
\end{grammarlens}

\subsection*{Your turn}
Say these aloud:

\begin{itemize}
\raggedright
  \item \emph{disculpe}
  \item \emph{disculpe}, then \emph{disculpa}, and hear which is the stranger's
  \item \emph{perdón}, then \emph{disculpe}
\end{itemize}

\begin{tcolorbox}[breakable,colback=teal!4,colframe=teal!35!black,title={Before you move on}]
What is inside \emph{disculpe}? (\emph{Culpa}, fault.) What does the \emph{dis-} do? (Undoes it.) Which form would you use with a stranger? (\emph{Disculpe}.)
\end{tcolorbox}
\section[con permiso]{con permiso — the one you say before, not after}

\label{lesson:ES-C283-permiso}



\begin{quote}
Before the new one: what does the \emph{dis-} do to \emph{culpa}? (Undoes it.)
\end{quote}

\subsection*{What to know first}
\begin{itemize}
\raggedright
  \item disculpe — the other \textquotedblleft{}excuse me\textquotedblright{}, for something done.
  \item por favor — the politeness you attach to a request.
\end{itemize}

\begin{sounds}
\begin{itemize}
\raggedright
  \item \emph{con per-MI-so}, stress on the \emph{mi}. The \emph{r} is one tap of the tongue, not the long English glide. $\to$ reference
\end{itemize}
\end{sounds}

\begin{cousinweb}
\textbf{con permiso} — \textquotedblleft{}excuse me,\textquotedblright{} said as you pass in front of someone.

\textbf{Permiso} is Latin \textbf{permissum}, from \emph{permittere}: \emph{per-}, through, plus \textbf{mittere}, to send. To permit is to let a thing through.

\emph{Mittere} was one of the most productive words Latin had, and English is full of it. A \textbf{mission} is a sending. A \textbf{missile} is a thing sent. To \textbf{dismiss} is to send away, and to \textbf{permit} is this very word, borrowed unchanged.

The two courtesies are not interchangeable, and the difference is time. \emph{Disculpe} is for something you have done. \emph{Con permiso} is for something you are about to do.
\end{cousinweb}

\begin{grammarlens}[title={What you've built}]
Politeness in Spanish is often a matter of tense rather than of warmth.

One formula looks backwards at a fault, the other forwards at an intrusion. Get the direction right and you will sound careful even when your grammar is not.
\end{grammarlens}

\subsection*{Your turn}
Say these aloud:

\begin{itemize}
\raggedright
  \item \emph{con permiso}
  \item \emph{disculpe}, then \emph{con permiso}, and place each one in time
  \item \emph{con permiso, por favor}
\end{itemize}

\begin{tcolorbox}[breakable,colback=teal!4,colframe=teal!35!black,title={Before you move on}]
What are the two pieces of \emph{permiso}? (\emph{Per}, through, and \emph{mittere}, to send.) Which courtesy comes before the act? (\emph{Con permiso}.)
\end{tcolorbox}
\section[amable]{amable — loveable, before it meant kind}

\label{lesson:ES-C283-amable}



\begin{quote}
Before the new one: which courtesy comes before the act? (\emph{Con permiso}.)
\end{quote}

\subsection*{What to know first}
\begin{itemize}
\raggedright
  \item gracias — what you say to someone who is this.
  \item disculpe — and what you say when you are not.
\end{itemize}

\begin{sounds}
\begin{itemize}
\raggedright
  \item \emph{a-MA-ble}, stress in the middle. The \emph{b} between vowels is soft — the lips barely close. $\to$ reference
\end{itemize}
\end{sounds}

\begin{cousinweb}
\textbf{amable} — \textquotedblleft{}kind.\textquotedblright{}

Latin \textbf{amabilis}, from \textbf{amare}, to love: loveable, a person it is easy to be fond of. Spanish moved the word from what you feel about the person to what the person does.

English has it twice over. \textbf{Amiable} came through French and kept the sense. \textbf{Amorous} kept the love, and an \textbf{amateur} is one who does a thing for love rather than for money — which is praise in the original and faint praise now.

The \textbf{-able} on the end is the same ending English writes on its own words. It arrived from the same Latin \emph{-abilis}, and it means then what it means now: able to have this done to it.
\end{cousinweb}

\begin{grammarlens}[title={What you've built}]
\emph{Muy amable} is the reply you have been missing.

You could thank someone; you could not say anything about them. One adjective closes that gap, and it costs you nothing new — \emph{muy} has been yours for a long while.
\end{grammarlens}

\subsection*{Your turn}
Say these aloud:

\begin{itemize}
\raggedright
  \item \emph{amable}, then \emph{muy amable}
  \item \emph{gracias, muy amable}
  \item \emph{amable} and \emph{amiable}, and hear the same word twice
\end{itemize}

\begin{tcolorbox}[breakable,colback=teal!4,colframe=teal!35!black,title={Before you move on}]
What did \emph{amabilis} mean in Latin? (Loveable.) Which Latin root is under it? (\emph{Amare}, to love.)
\end{tcolorbox}
\section[mil gracias]{mil gracias — a number that stopped counting}

\label{lesson:ES-C283-mil-gracias}



\begin{quote}
Before the new one: what did \emph{amabilis} mean? (Loveable.)
\end{quote}

\subsection*{What to know first}
\begin{itemize}
\raggedright
  \item gracias — the half of this you already have.
  \item amable — what you say about the person you are thanking.
\end{itemize}

\begin{sounds}
\begin{itemize}
\raggedright
  \item \emph{mil GRA-cias}. The \emph{l} stays clear and forward, never darkening the way an English \emph{l} does at the end of a word. $\to$ reference
\end{itemize}
\end{sounds}

\begin{cousinweb}
\textbf{mil gracias} — \textquotedblleft{}a thousand thanks.\textquotedblright{}

\emph{Gracias} has been yours since near the beginning. \textbf{Mil} is Latin \textbf{mille}, a thousand, and it is the same \emph{mille} buried in English \textbf{mile}: the Roman \textbf{mille passuum}, a thousand paces, which is why a mile is the awkward length it is and not a round number of anything.

Notice what the phrase is doing. Spanish is not measuring gratitude here; it is reaching for a large number and stopping. Every language keeps a few numbers that have given up counting — English says thanks a million and means the same nothing-in-particular.

\textbf{Millimetre}, \textbf{million} and \textbf{millennium} are the honest arithmetic uses of \emph{mille}, and they arrived from the book rather than from the mouth.
\end{cousinweb}

\begin{grammarlens}[title={What you've built}]
Four courtesies, and not one of them is a new idea.

Each is a word you had, or a number you had, put to a use you did not know about. Vocabulary grows fastest at the joins.
\end{grammarlens}

\subsection*{Your turn}
Say these aloud:

\begin{itemize}
\raggedright
  \item \emph{mil gracias}
  \item \emph{gracias}, then \emph{mil gracias}, and hear the weight change
  \item \emph{mil gracias, muy amable}
\end{itemize}

\begin{tcolorbox}[breakable,colback=teal!4,colframe=teal!35!black,title={Before you move on}]
What is \emph{mil} in Latin? (\emph{Mille}, a thousand.) What has a mile got to do with it? (It was a thousand paces.)
\end{tcolorbox}
\section[no hay de qué]{no hay de qué — four words you already own}

\label{lesson:ES-C283-no-hay-de-que}



\begin{quote}
Before the last of the five: what is \emph{mil} in Latin? (\emph{Mille}, a thousand.) And what did \emph{amabilis} mean? (Loveable.)
\end{quote}

\subsection*{What to know first}
\begin{itemize}
\raggedright
  \item de nada — the other reply to thanks, and the same idea.
  \item mil gracias — what this answers.
\end{itemize}

\begin{sounds}
\begin{itemize}
\raggedright
  \item \emph{no ay de KE}, four quick syllables run together. The written accent on \emph{qué} marks the question word, not a change of stress. $\to$ reference
\end{itemize}
\end{sounds}

\begin{cousinweb}
\textbf{no hay de qué} — \textquotedblleft{}not at all,\textquotedblright{} the reply when someone thanks you.

Every piece is already yours. \textbf{No}, \textbf{hay}, \textbf{de}, \textbf{qué} — nothing new in the list. What is new is the assembly: literally \textquotedblleft{}there is nothing of which,\textquotedblright{} with the thing not worth mentioning left out of the sentence entirely.

Spanish already gave you \textbf{de nada}, \textquotedblleft{}of nothing,\textquotedblright{} for the same moment. The two are one idea built twice: whatever you are thanking me for, there was nothing in it.

That is worth more than a fifth word. A courtesy formula is rarely new vocabulary. It is old vocabulary frozen in an order nobody parses any more — and once you have taken one apart, the rest stop being phrases to memorise.
\end{cousinweb}

\begin{grammarlens}[title={What you've built}]
Five ways to be polite: \emph{disculpe}, \emph{con permiso}, \emph{amable}, \emph{mil gracias}, \emph{no hay de qué}.

Two of them look backwards and two look forwards; the fifth is made of scraps. And the fifth is the one that shows you what the others are — sentences that stopped being sentences.
\end{grammarlens}

\subsection*{Your turn}
Say these aloud:

\begin{itemize}
\raggedright
  \item \emph{no hay de qué}
  \item its four pieces slowly — \emph{no}, \emph{hay}, \emph{de}, \emph{qué}
  \item all five in order — \emph{disculpe}, \emph{con permiso}, \emph{amable}, \emph{mil gracias}, \emph{no hay de qué}
\end{itemize}

\begin{tcolorbox}[breakable,colback=teal!4,colframe=teal!35!black,title={Before you move on}]
How many new words are in \emph{no hay de qué}? (None — only the assembly is new.) Which reply to thanks did you already have? (\emph{De nada}.) And what is inside \emph{disculpe}? (\emph{Culpa}, fault.)
\end{tcolorbox}
